O Pugh compara alternativas a um datum, no mesmo nível de abstração, e orienta a discussão sem transformar o total em decisão absoluta \citep{mitpugh2015}. Nesta primeira seleção, a equipe avalia diferenças entre estruturas funcionais. A confirmação de aceleração, força, volume, vida e custo pertence à comparação das concepções, conforme as soluções e análises disponíveis. Os 23 requisitos, suas metas e seus protocolos permanecem obrigatórios.

A comparação atual usa M2 como datum e a escala qualitativa $-1$, $0$, $+1$, sem ponderação adicional. $+1$ indica vantagem funcional identificável; $-1$, desvantagem identificável; $0$, a mesma cobertura ou o mesmo aspecto descrito. Ausência de base recebe AV, sem ser convertida em zero. Os pesos derivados do Mudge/Pareto e a escala de cinco níveis permanecem no histórico e na segunda matriz. Não foram inventados pesos para os novos critérios funcionais, nem usados pontos do QFD como se fossem preferências adicionais.

\begin{table}[H]
\centering\footnotesize
\caption{Primeira Pugh no nível das estruturas funcionais}\label{tab:pugh-funcional}
\begin{tabularx}{\textwidth}{L{0.9cm}Y*{3}{>{\centering\arraybackslash}p{1.05cm}}}
\toprule\textbf{ID}&\textbf{Aspecto comparado}&\textbf{M1}&\textbf{M2}&\textbf{M3}\\\midrule
CF01&Cobertura de proteção, retenção e movimentação&0&0&0\\
CF02&Barreiras pré-existentes no caminho de inspeção&+1&0&0\\
CF03&Continuidade do acondicionamento no distribuidor&$-1$&0&0\\
CF04&Cobertura da informação obrigatória ao usuário&0&0&0\\
CF05&Cobertura de retorno, manutenção e fim de vida&0&0&0\\\midrule
&Vantagens / desvantagens / equivalências&1/1/3&0/0/5&0/0/5\\
&\textbf{Saldo qualitativo}&\textbf{0}&\textbf{0}&\textbf{0}\\
\bottomrule
\end{tabularx}
\fonte{Julgamentos da equipe a partir das fronteiras, dos fluxos e do catálogo funcional. A planilha registra justificativa e origem por alternativa.}
\end{table}

CF01, CF04 e CF05 receberam zero porque as três estruturas incluem os mesmos efeitos obrigatórios de proteção, informação e suporte ao reuso. Esses zeros não afirmam igualdade de desempenho físico. OP-01 não melhora automaticamente esses efeitos e continua fora da configuração base.

Em CF02, M1 elimina a embalagem de origem do caminho de observação da cerâmica após o preparo no distribuidor. A barreira do novo dispositivo permanece. A nota $+1$ registra uma facilidade potencial de acesso, sem demonstrar a área crítica exigida por RP4. M2 e M3 mantêm a mesma barreira de origem; registrar choque não a torna transparente. Em CF03, retirar a caixa e transferir a peça acrescenta uma operação em M1. A nota $-1$ representa essa intervenção necessária, sem atribuir tempo ou taxa de avaria ainda não medidos.

\noindent\textbf{Decisão de planejamento.} Em 15 de outubro de 2026, a equipe priorizou M1: retirar a embalagem original no distribuidor e desenvolver proteção direta da peça. A comparação termina em empate, e a preferência explícita orienta estudar o acesso mais direto, assumindo a contrapartida do preparo. M1 não foi declarada vencedora numérica. M2 permanece alternativa de contingência; M3 não é adotada na base.

A escolha encerra a definição da direção funcional para prosseguir, condicionada à viabilidade da proteção direta. O preparo deve incluir retirada controlada, conferência inicial, transferência segura, identidade da unidade e destino da embalagem. Definir pressão e área de contato, inspeção com o dispositivo fechado e contenção na abertura. Se essas condições não puderem ser resolvidas, reabrir a fronteira funcional.

A planilha conserva os 23 critérios do cliente e explicita seu encaminhamento à segunda Pugh, onde serão comparadas C1/C2/C3 sob M1. A direção funcional apoia-se na descrição comparativa; os protocolos permanecem como base para analisar e confirmar o desempenho das soluções físicas. A introdução posterior de OP-01 na M1 exigiria revisão de escopo e benefício, mantendo a fronteira direta da peça; isso não corresponde à M3, definida sobre o volume embalado.
